A autonomia artificial exige repensar o papel da recompensa no Aprendizado por Reforço. Em abordagens tradicionais, o agente aprende a partir de recompensas externas definidas pelo projetista, o que pode produzir comportamentos eficientes, mas limita a emergência de formas mais autônomas de ação. Inspirado pela Inteligência Artificial Enativa, este trabalho investiga agentes cuja ação é orientada por condições internas de viabilidade. Para isso, utiliza-se a tarefa experimental \textit{Health-Gathering}, da plataforma VizDoom, na qual um agente perde vida continuamente em uma sala com chão ácido e percebe o ambiente por meio da visão. Em vez de recompensar explicitamente a coleta de \textit{medikits}, itens que restauram a vida do agente, introduz-se uma forma artificial de precariedade corporal. À medida que o agente permanece sem coletar \textit{medikits}, sua visão é degradada, sendo restaurada quando ele os coleta. A motivação do agente é produzida por uma recompensa intrínseca baseada em curiosidade, calculada por meio do método \textit{Random Network Distillation} (RND). Como a curiosidade depende das observações visuais, a degradação perceptiva reduz a capacidade do agente de obter recompensa intrínseca. Assim, manter-se viável torna-se condição para continuar percebendo, explorando e agindo. A política aprendida será analisada com técnicas de interpretabilidade visual, buscando compreender como precariedade, percepção e curiosidade podem favorecer comportamentos adaptativos sem impor diretamente uma solução por recompensas extrínsecas. Sendo assim, se a autonomia nasce da precariedade, então ser inteligente é, fundamentalmente, estar engajado em um processo contínuo de negociação com a possibilidade de deixar de existir.

% Separe as palavras-chave por ponto
\palavraschave{Agentes autônomos. Inteligência Artificial Enativa. Curiosidade. Aprendizado por Reforço Profundo. Precariedade. Motivação intrínseca}
